\section*{第 \ref{ch:b2:pt-paths} 章——地质温压计与 P--T--t 轨迹}

\begin{solution}{exo:b2:pt-paths:1}
(a) A 绿泥石带，B 黑云母带，C 石榴子石带，D 十字石带，E 蓝晶石带，
F 夕线石带。(b) 共五条，即黑云母、石榴子石、十字石、蓝晶石和夕线石
等变质线，各位于相邻两个样品之间，变质级自 A 向 F 递增。(c) 等变质线
只说明某个反应发生过；\omterm{def:b2:pt-paths:barrovian}{变质场梯度}需要压强和温度，
这些数据来自经过标定的反应，或来自对
每块岩石峰期\omterm{def:b2:metamorphic-facies:assemblage}{矿物组合}进行的温压计计算。
\end{solution}

\begin{solution}{exo:b2:pt-paths:2}
(a) $(X_{\text{Mg}}/X_{\text{Fe}})_{\text{grt}} = 0.10/0.80 = 0.125$，
$(X_{\text{Mg}}/X_{\text{Fe}})_{\text{bt}} = 0.45/0.50 = 0.900$；$K_D = 0.1389$，$\ln K_D =
-1.9741$。(b) $3R = 24.942$；分母 $20.19 + 24.942\times1.9741 = 69.428$；在
\qty{5}{kbar} 下，$T = 54185/69.428 = \qty{780.5}{K} = \qty{507}{\celsius}$；在
\qty{7}{kbar} 下，$T = (52630 + 2177)/69.428 = \qty{789.4}{K} = \qty{516}{\celsius}$。
(c) 两千巴的压强变化仅使其改变 \qty{9}{K}：它是一支温度计。
\end{solution}

\begin{solution}{exo:b2:pt-paths:3}
(a) $z = P/(\bar\rho g) = 6.0\times10^8/(2800\times9.81) = \qty{21.8}{km}$。
(b) $(580 - 10)/21.8 = \qty{26}{\celsius\per\km}$。(c) 在 \qty{580}{\celsius} 时，该线
位于 $3.67 + 0.8\times2.10 = \qty{5.35}{kbar}$ 处；在 \qty{6.0}{kbar} 时，岩石位于
高压一侧：蓝晶石。
\end{solution}

\begin{solution}{exo:b2:pt-paths:4}
(a) 反应物 $2\times140.693 + 3\times18.015 = \qty{335.431}{g}$；生成物 $277.112 +
58.320 = \qty{335.432}{g}$：在取舍误差范围内配平。(b) $2\times43.65 =
\qty{87.30}{cm^3}$ 变为 $107.50 + 24.63 = \qty{132.13}{cm^3}$：$+\qty{51.4}{\percent}$。
(c) 橄榄石 $281.386/87.30 = \qty{3.223}{g.cm^{-3}}$；生成物 $335.432/132.13 =
\qty{2.539}{g.cm^{-3}}$。(d) 岩石体积膨胀约一半，密度变小；它
还将水储存在其矿物之中。
\end{solution}

\begin{solution}{exo:b2:pt-paths:5}
(a) 均值在 \qty{1072.15}{K} 时为 $-1.206$，在 \qty{823.15}{K} 时为 $-1.782$。取 $1/T =
9.3271\times10^{-4}$ 和 $1.21485\times10^{-3}$，$a = (-1.206 + 1.782)/(9.3271\times10^{-4}
- 1.21485\times10^{-3}) = -2041.5$ K，且 $b = -1.206 + 2041.5/1072.15 = 0.6982$。
(b) $\Delta_r H + 2070\times0.311 = 2041.5\times24.942 = \qty{50920}{J}$，故 $\Delta_r H =
\qty{50276}{J}$；$\Delta_r S = 0.6982\times24.942 = \qty{17.41}{J.K^{-1}}$。(c) 在
\qty{5}{kbar} 下：对于 $K_D = 0.18$，得 588.1，对比 \qty{587.5}{\celsius}；对于 $0.10$，得 419
对比 \qty{425}{\celsius}；对于 $0.30$，得 819，对比 \qty{806}{\celsius}。在实验范围内
两种标定结果一致，尽管其常数分别相差
\qty{2.4}{kJ} 和 \qty{2.8}{J.K^{-1}}；$\Delta_r H$ 与 $\Delta_r S$ 的误差
在此范围内相互抵消，在范围之外则增大。标定结果不宜用于远离
其实验条件的场合。
\end{solution}

\begin{solution}{exo:b2:pt-paths:6}
(a) \omterm{def:b2:pt-paths:thermobarometry}{交换反应}：$(113.12 + 154.30) - (115.32 + 149.65) = \qty{+2.45}{cm^3} =
\qty{0.245}{J.bar^{-1}}$，对比 0.311。GASP 反应：$125.28 + 2\times44.15 + 22.69 -
3\times100.79 = \qty{-66.10}{cm^3} = \qty{-6.61}{J.bar^{-1}}$，对比 $-6.21$，即
在 \qty{800}{\celsius} 和 \qty{10}{kbar} 下的标定值，此时热膨胀与
压缩改变了体积。两种情形下，GASP 的变化量为
\omterm{def:b2:pt-paths:thermobarometry}{交换反应}的 20 至 30 倍。(b) $\partial T/\partial P = 0.311/62.96 = \qty{4.9e-3}{K/bar}$：
\qty{4.9}{K}，即每千巴。(c) 该线的斜率为 $22.80 + 8.314\times3\ln 0.20/6.21
= 22.80 - 6.46 = \qty{16.3}{bar/K}$；\qty{17}{K} 的变化使压强移动
\qty{280}{bar}。
\end{solution}

\begin{solution}{exo:b2:pt-paths:7}
(a) $P_0 = 22.80\times600 - 1093 = \qty{12587}{bar}$；$3\ln(0.08/0.30) = -3.9653$；$RT =
8.314\times873.15 = \qty{7259.4}{J}$；$P = 12587 + 7259.4\times(-3.9653)/6.21 = 12587 -
4635 = \qty{7952}{bar}$。(b) $3\ln(0.08/0.35) = -4.4277$，$P = 12587 - 5176 =
\qty{7411}{bar}$：少 \qty{0.54}{kbar}，因此斜长石必须谨慎分析。
(c) $7952\times10^5/(2800\times9.81) = \qty{29.0}{km}$。
\end{solution}

\begin{solution}{exo:b2:pt-paths:8}
(a) $z = P/(\bar\rho g)$：\qty{118.9}{km}、\qty{34.0}{km} 和 \qty{15.3}{km}。
(b) $(118.9 - 34.0)/2.2 = \qty{38.6}{km/Ma} = \qty{3.9}{cm.\text{年}^{-1}}$ 和 $(34.0 -
15.3)/1.1 = \qty{17.0}{km/Ma} = \qty{1.7}{cm.\text{年}^{-1}}$，略高于已发表的
3.4 和 1.6：他们将压强换算为深度时，所得的深度
差略小。(c) $\sigma = \sqrt{0.9^2 + 0.9^2} = \qty{1.3}{Ma}$，因此 $2.2 \pm 1.3$ 百万年，
速率在 $84.9/3.47 = 24$ 与 $84.9/0.93 = \qty{91}{km/Ma}$ 之间，
即 2.4 至 \qty{9.1}{cm.\text{年}^{-1}}：无论如何都很快。(d) 即使最低速率，也约为
模型中 \qty{0.5}{mm.\text{年}^{-1}} 的五十倍：仅靠侵蚀无法做到。
\end{solution}

\begin{solution}{exo:b2:pt-paths:9}
(a) $\partial T/\partial\ln K_D = 3RT/(\Delta_r S - 3R\ln K_D) = 24.942\times860.6/62.96 =
\qty{341}{K}$；$\partial T/\partial P = 0.311/62.96 = \qty{4.9}{K}$，即每千巴。
(b) $\sigma_T = \sqrt{(341\times0.05)^2 + 4.9^2} = \sqrt{17.0^2 + 4.9^2} =
\qty{17.7}{K}$。(c) $K_D$ 的分析误差起主导作用，仅此一项即占
最佳标定结果中 \qty{25}{\celsius} 的三分之二，这些标定还包括
标定本身的误差。
\end{solution}

\begin{solution}{exo:b2:pt-paths:10}
(a) 石榴子石 $0.10/0.75 = 0.1333$，黑云母 $0.42/0.52 = 0.8077$，$K_D = 0.1650$，$\ln K_D =
-1.8018$，分母 $20.19 + 24.942\times1.8018 = 65.130$；$3\ln(0.07/0.25) = -3.8189$。
在 \qty{6000}{bar} 下：$T = 54496/65.130 = \qty{836.7}{K} = \qty{563.6}{\celsius}$，然后 $P
= 22.80\times563.6 - 1093 + 8.314\times836.7\times(-3.8189)/6.21 = \qty{7479}{bar}$。
在 \qty{7479}{bar} 下：$T = \qty{570.6}{\celsius}$，$P = \qty{7603}{bar}$。在
\qty{7603}{bar} 下：$T = \qty{571.2}{\celsius}$，$P = \qty{7614}{bar}$；再算一次，$T =
\qty{571.3}{\celsius}$。岩石在约 \qty{571}{\celsius} 和
\qty{7.6}{kbar} 下达到平衡。(b) 在 \qty{571}{\celsius} 时，蓝晶石 $=$ 夕线石的反应线位于 $3.67
+ 0.71\times2.10 = \qty{5.16}{kbar}$：蓝晶石。(c) 该压力计是针对
蓝晶石标定的；若用夕线石，则反应及其反应线不同，因为这两种同质多象变体
的吉布斯自由能并不相同。
\end{solution}

\begin{solution}{exo:b2:pt-paths:11}
(a) $\tau = L^2/(\pi^2\kappa)$：对于 \qty{35}{km}，$1.225\times10^9/(9.870\times10^{-6}) =
\qty{1.24e14}{s} = \qty{3.9}{Ma}$；对于 \qty{70}{km}，时间长四倍，
\qty{15.7}{Ma}。(b) 7000 万年是这些时间的 4 至 18 倍：热量
通过增厚的地壳迁移，远快于侵蚀将其去除。(c) 岩石
一次性到达最大深度，在叠加作用中完成；热量在随后的
数千万年中才到达，而侵蚀正在缓慢将其抬升：温度仍继续
在减压过程中升高一段时间。(d) 若侵蚀在 700 万年内
就将岩石抬升，此时传导尚未使其升温：热峰值更低，
出现较早、压强较高，并沿一条冷路径返回地表。
\end{solution}

\begin{solution}{exo:b2:pt-paths:12}
(a) 地壳的重量为 $35000\times2800\times9.81 = \qty{0.961}{GPa}$；其余
\qty{1.839}{GPa} 需要 $1.839\times10^9/(3300\times9.81) = \qty{56.8}{km}$ 厚的地幔，深度
为 \qty{91.8}{km}，与原研究的 \qty{90}{km} 或更大值相符。(b) $(4.0 -
0.961)\times10^9/(3300\times9.81) = \qty{93.9}{km}$，深度 \qty{128.9}{km}。
(c) $700/91.8 = \qty{7.6}{\celsius\per\km}$：压力与
温度之比很高，属于俯冲带系列。(d) $(91.8 - 35)/34 = \qty{1.7}{Ma}$。
\end{solution}

\begin{solution}{pb:b2:pt-paths:1}
\noindent\textbf{第一部分。}

\textbf{1.} 石榴子石核心 $0.080/0.760 = 0.1053$，黑云母 $0.40/0.58 = 0.6897$；$K_D =
0.1527$，$\ln K_D = -1.8793$。

\textbf{2.} 边缘 $0.120/0.720 = 0.1667$，基质黑云母 $0.46/0.53 = 0.8679$；$K_D =
0.1921$，$\ln K_D = -1.6497$。

\textbf{3.} $3R = 24.942$；分子 $52630 + 0.311\times6000 = 54496$；分母 $20.19
+ 24.942\times1.8793 = 67.064$；$T = \qty{812.6}{K} = \qty{539.5}{\celsius}$。

\textbf{4.} 分母 $20.19 + 24.942\times1.6497 = 61.337$；$T = 54496/61.337 =
\qty{888.5}{K} = \qty{615.3}{\celsius}$。

\textbf{5.} 对核心为 $0.311\times1000/67.064 = \qty{4.6}{K}$，对边缘为 $0.311\times1000/61.337
= \qty{5.1}{K}$。

\textbf{6.} \omterm{def:b2:pt-paths:thermobarometry}{地质温度计}只适用于曾经处于平衡状态的矿物。包含在核心中的
黑云母在核心生长时便与岩石其余部分隔绝；
基质黑云母最后一次与边缘交换。将核心与基质配对，会把
两个不同时刻混在一起。

\noindent\textbf{第二部分。}

\textbf{7.} 核心 $22.80\times539.5 - 1093 = \qty{11208}{bar}$；边缘 $22.80\times615.3 -
1093 = \qty{12936}{bar}$。

\textbf{8.} $3\ln(0.050/0.25) = -4.8283$；$RT = 8.314\times812.65 = \qty{6756.4}{J}$；$P =
11208 - 6756.4\times4.8283/6.21 = 11208 - 5253 = \qty{5955}{bar}$。

\textbf{9.} $3\ln(0.060/0.28) = -4.6213$；$RT = 8.314\times888.45 = \qty{7386.6}{J}$；$P
= 12936 - 5497 = \qty{7439}{bar}$。

\textbf{10.} 当 \qty{7439}{bar} 时：$T = (52630 + 2314)/61.337 = \qty{895.8}{K} =
\qty{622.6}{\celsius}$；在该温度下 $P = 13102 - 5542 = \qty{7560}{bar}$；然后 $T
= \qty{623.2}{\celsius}$。边缘大约在 \qty{623}{\celsius} 和 \qty{7.56}{kbar} 下生长。

\textbf{11.} 当 \qty{5954}{bar} 时，$T = \qty{539.2}{\celsius}$，且 $P = \qty{5950}{bar}$：
核心大约在 \qty{539}{\celsius} 和 \qty{5.95}{kbar} 下生长。最初的估计
\qty{6}{kbar} 已经很接近，且温度计每千巴只移动 \qty{4.6}{K}。

\textbf{12.} 核心 $5950\times10^5/(2800\times9.81) = \qty{21.7}{km}$；边缘
$7560\times10^5/27468 = \qty{27.5}{km}$。

\textbf{13.} 在 \qty{539}{\celsius} 时，该线位于 $3.67 + 0.39\times2.10 =
\qty{4.49}{kbar}$，低于 \qty{5.95}{kbar}：为蓝晶石，与包含在
核心中的蓝晶石一致。在 \qty{623}{\celsius} 时，该线位于 $5.77 + 0.23\times2.09 = \qty{6.25}{kbar}$，
低于 \qty{7.56}{kbar}：仍为蓝晶石。

\noindent\textbf{第三部分。}

\textbf{14.} $+\qty{84}{\celsius}$ 与 $+\qty{1.6}{kbar}$（约深 \qty{6}{km}）：岩石
正在被埋藏并加热，处于\omterm{def:b2:pt-paths:paths}{进变质轨迹}上。

\textbf{15.} 在 \qty{623}{\celsius} 或更高温度、低于 \qty{6.25}{kbar} 的条件下：边缘生长之后，岩石
至少损失了 \qty{1.3}{kbar}，约相当于 \qty{5}{km} 的上覆岩层，同时温度保持至少
同样高。它是在仍然炽热时被\omterm{def:b2:pt-paths:exhumation}{折返}的。

\textbf{16.} 绿泥石是低级含水矿物，它记录了晚期\omterm{def:b2:pt-paths:paths}{退变质
轨迹}，即在浅部、有水条件下冷却。

\textbf{17.} 埋藏并加热至边缘生长，随后减压但不降温，再冷却：
若以温度为横轴向右、压强向上作图，轨迹呈顺时针方向。压强峰值
最先出现在边缘；热峰值与之同时或更晚，出现在
减压过程中。

\textbf{18.} 是的：该模型同样先达到压强峰值，其热峰值
为 \qty{593}{\celsius}、\qty{7.8}{kbar}，出现在上升过程中，接近片岩
边缘的条件。增厚之后再缓慢剥蚀，可以解释该片岩的轨迹。

\noindent\textbf{第四部分。}

\textbf{19.} $(623 - 405)/(412.0 - 398.0) = 218/14.0 = \qty{15.6}{\celsius}$ 每
百万年。

\textbf{20.} 岩石的冷却速度约为假设的 1.5 倍。一种矿物
冷却更快，逸失氩的时间更少，因此在更高温度下封闭：
真实的封闭温度略高于 \qty{405}{\celsius}，相应的封闭
深度也略深。

\textbf{21.} $(405 - 10)/30 = \qty{13.2}{km}$。

\textbf{22.} $27.5 - 13.2 = \qty{14.3}{km}$（若采用未取整的深度 27.52 和 13.17，
则为 \qty{14.35}{km}）。

\textbf{23.} $\sqrt{1.5^2 + 1.0^2} = \qty{1.8}{Ma}$：$14.0 \pm \qty{1.8}{Ma}$，因此速率
介于 $14.35/15.8 = 0.91$ 与 $14.35/12.2 = \qty{1.18}{km/Ma}$ 之间。

\textbf{24.} $14.35/14.0 = \qty{1.03}{km}$ 每百万年，即
\textbf{\qty[text-series-to-math=true]{1.0}{mm.\text{年}^{-1}}}，约慢三十倍，比阿尔卑斯
最深岩石\omterm{def:b2:pt-paths:exhumation}{折返}的第一阶段，并约为增厚模型侵蚀速率的
两倍。
\end{solution}
